\section{Introduction}
\subsection{Purpose and scope}
This report establishes an illustrative geotechnical design basis for a reinforced-concrete
mid-rise building with one basement. It demonstrates the evaluation of a ground model,
foundation alternatives, axial pile resistance, settlement, groundwater risks and construction
quality controls. The assessment covers:
\begin{itemize}
  \item interpretation of a representative three-borehole investigation to
        \SI{25}{\metre} depth;
  \item selection of characteristic engineering parameters;
  \item screening of pads, raft and deep-foundation options;
  \item preliminary geotechnical sizing of bored piles; and
  \item requirements for verification, monitoring and quality assurance.
\end{itemize}

This report does not design retaining walls, permanent basement waterproofing, pile caps,
structural reinforcement or temporary works. Those items require coordinated design after the
project geometry and actions are confirmed.

\subsection{Assumed development}
\begin{table}[H]
\centering
\caption{Assumed project design basis.}
\label{tab:project-basis}
\begin{tabularx}{\textwidth}{@{}L{47mm}Y@{}}
\toprule
\rowcolor{navy!8}
\textbf{Item} & \textbf{Adopted basis} \\
\midrule
Building & Eight-storey reinforced-concrete frame with one basement \\
Plan dimensions & Approximately \SI{22}{\metre} by \SI{30}{\metre} \\
Basement / pile-cap level & About \SI{2.7}{\metre} below existing grade \\
Service column reactions & \SIrange{700}{1800}{\kilo\newton}; core reaction about \SI{9000}{\kilo\newton} \\
Settlement criteria & \SI{40}{\milli\metre} total; angular distortion preferably less than 1:500 \\
Adjacent assets & Low-rise buildings and buried utilities near the property boundary \\
Design life & 50 years, subject to confirmation by the owner \\
\bottomrule
\end{tabularx}
\end{table}

\subsection{Information hierarchy}
Where conflicts arise, project-specific survey, certified field logs and laboratory results take
precedence over this illustrative basis. Design should follow the governing building code and
the contract-adopted editions of recognized geotechnical standards. Parameter selection should
use engineering judgment and load-test back-analysis rather than uncorrected SPT values alone
\cite{peck1974,bowles1996}.

% =============================================================================
